%% ============================================================================
%%  Drosophila-Lang: un lenguaje esotérico Turing-completo sobre el conectoma
%%  completo de Drosophila melanogaster
%%
%%  Compilación:  pdflatex drosophila_lang.tex  (dos pasadas)
%%  Documento autocontenido: sin figuras externas ni ficheros .bib.
%% ============================================================================
\documentclass[11pt,a4paper]{article}

%% --- Codificación, idioma y tipografía --------------------------------------
\usepackage[utf8]{inputenc}
\usepackage[T1]{fontenc}
\usepackage[english,spanish,es-noshorthands,es-tabla]{babel}
\usepackage{lmodern}
\usepackage{microtype}
\usepackage[margin=2.5cm]{geometry}

%% --- Matemáticas ------------------------------------------------------------
\usepackage{amsmath,amssymb,amsthm,mathtools}

%% --- Figuras, tablas y gráficos ---------------------------------------------
\usepackage{graphicx}
\usepackage{booktabs}
\usepackage{tabularx}
\usepackage{multirow}
\usepackage{array}
\usepackage{xcolor}
\usepackage{tikz}
\usetikzlibrary{arrows.meta,positioning,calc,fit,backgrounds,shapes.geometric}
\usepackage{pgfplots}
\pgfplotsset{compat=1.16}
\usepackage[font=small,labelfont=bf]{caption}
\usepackage{subcaption}

%% --- Código y algoritmos ----------------------------------------------------
\usepackage{listings}
\usepackage{algorithm}
\usepackage{algpseudocode}
\usepackage{enumitem}

%% --- Referencias cruzadas (hyperref antes de cleveref) ----------------------
\usepackage[colorlinks=true,linkcolor=blue!60!black,citecolor=green!40!black,
            urlcolor=blue!60!black]{hyperref}
\usepackage[spanish,nameinlink,noabbrev]{cleveref}

\hypersetup{
  pdftitle={Drosophila-Lang: un lenguaje esotérico Turing-completo sobre el
            conectoma completo de Drosophila melanogaster},
  pdfauthor={Nombre Apellido},
  pdfkeywords={conectoma, FlyWire, Drosophila, lenguajes esotéricos,
               completitud de Turing, cuerpo en champiñón, plasticidad dopaminérgica}
}

%% --- Entornos de teoremas ----------------------------------------------------
\theoremstyle{plain}
\newtheorem{teorema}{Teorema}
\newtheorem{lema}{Lema}
\newtheorem{proposicion}{Proposición}
\newtheorem{corolario}{Corolario}
\theoremstyle{definition}
\newtheorem{definicion}{Definición}
\theoremstyle{remark}
\newtheorem{observacion}{Observación}

\crefname{teorema}{teorema}{teoremas}
\crefname{lema}{lema}{lemas}
\crefname{proposicion}{proposición}{proposiciones}
\crefname{corolario}{corolario}{corolarios}
\crefname{definicion}{definición}{definiciones}
\crefname{algorithm}{algoritmo}{algoritmos}
\crefname{lstlisting}{listado}{listados}
\crefname{table}{tabla}{tablas}
\Crefname{table}{Tabla}{Tablas}

%% --- Algoritmos en español ---------------------------------------------------
\floatname{algorithm}{Algoritmo}
\renewcommand{\listalgorithmname}{Lista de algoritmos}
\renewcommand{\algorithmicrequire}{\textbf{Entrada:}}
\renewcommand{\algorithmicensure}{\textbf{Salida:}}
\algrenewcommand\algorithmicif{\textbf{si}}
\algrenewcommand\algorithmicthen{\textbf{entonces}}
\algrenewcommand\algorithmicelse{\textbf{si no}}
\algrenewcommand\algorithmicfor{\textbf{para}}
\algrenewcommand\algorithmicdo{\textbf{hacer}}
\algrenewcommand\algorithmicend{\textbf{fin}}
\algrenewcommand\algorithmicreturn{\textbf{devolver}}
\algrenewcommand\algorithmicwhile{\textbf{mientras}}

%% --- Listados de código ------------------------------------------------------
\definecolor{fondocodigo}{RGB}{247,247,244}
\definecolor{comentario}{RGB}{110,120,110}
\definecolor{cadena}{RGB}{160,60,20}
\definecolor{clave}{RGB}{20,60,140}
\lstset{
  basicstyle=\ttfamily\footnotesize,
  backgroundcolor=\color{fondocodigo},
  frame=single, rulecolor=\color{black!20},
  columns=fullflexible, keepspaces=true,
  breaklines=true, breakatwhitespace=true,
  showstringspaces=false, upquote=true,
  xleftmargin=0.5em, xrightmargin=0.5em, aboveskip=0.8em, belowskip=0.8em,
  literate={á}{{\'a}}1 {é}{{\'e}}1 {í}{{\'i}}1 {ó}{{\'o}}1 {ú}{{\'u}}1
           {Á}{{\'A}}1 {É}{{\'E}}1 {Í}{{\'I}}1 {Ó}{{\'O}}1 {Ú}{{\'U}}1
           {ñ}{{\~n}}1 {Ñ}{{\~N}}1 {¡}{{!`}}1 {¿}{{?`}}1 {ü}{{\"u}}1
           {§}{{\S}}1 {≈}{{$\approx$}}1 {⇝}{{$\rightsquigarrow$}}1 {†}{{\dag}}1
}
\lstdefinelanguage{Drosophila}{
  morecomment=[l]{;;},
  morestring=[b]",
  morekeywords={x_x},
  sensitive=true
}
\lstdefinestyle{dros}{
  language=Drosophila,
  commentstyle=\color{comentario}\itshape,
  stringstyle=\color{cadena},
  keywordstyle=\color{clave}\bfseries
}
\lstdefinestyle{py}{
  language=Python,
  commentstyle=\color{comentario}\itshape,
  stringstyle=\color{cadena},
  keywordstyle=\color{clave}\bfseries
}
\renewcommand{\lstlistingname}{Listado}

%% --- Macros ------------------------------------------------------------------
\newcommand{\dros}{\textup{\textsc{Drosophila-Lang}}}
\newcommand{\Dm}{\textit{Drosophila melanogaster}}
\newcommand{\PAM}{\mathrm{PAM}}
\newcommand{\PPL}{\mathrm{PPL1}}
\newcommand{\GF}{\mathrm{GF}}
\newcommand{\KC}{\mathrm{KC}}
\newcommand{\MBON}{\mathrm{MBON}}
\newcommand{\ORN}{\mathrm{ORN}}
\newcommand{\sig}{\sigma}
\newcommand{\ind}{\mathbf{1}}
\newcommand{\N}{\mathbb{N}}
\newcommand{\Z}{\mathbb{Z}}
\newcommand{\R}{\mathbb{R}}

%% ============================================================================
\title{\textbf{Drosophila-Lang}: un lenguaje de programación esotérico y
       Turing-completo ejecutado sobre el conectoma completo de
       \textit{Drosophila melanogaster}}

\author{
  Nombre Apellido\thanks{Autor de correspondencia: \texttt{correo@institucion.org}}\\[2pt]
  \small Departamento, Institución, Ciudad, País
}

\date{25 de septiembre de 2026}

\begin{document}
\maketitle

%% ============================================================================
\begin{abstract}
Presentamos \dros{}, un lenguaje de programación esotérico cuyo modelo de
ejecución no es una máquina de von Neumann, sino el cerebro de una mosca de
la fruta. Los programas se ejecutan sobre el conectoma adulto completo FlyWire
v783 (138\,639 neuronas, 2\,700\,513 conexiones con al menos cinco sinapsis y
34\,153\,566 sinapsis en total), con signo sináptico derivado del
neurotransmisor predicho. El código fuente se escribe como un mapa de
glomérulos del lóbulo antenal: cada instrucción se asigna, mediante un
procedimiento de reclutamiento guiado por la anatomía, a una neurona real
identificada (neuronas de proyección, células de Kenyon, neuronas de salida
del cuerpo en champiñón, dopaminérgicas PAM y PPL1, neuronas locales del
cuerno lateral, receptores olfativos, motoneuronas, la descendente DNa02, la
motoneurona MN9 y la Fibra Gigante DNp01). La memoria del programa no reside
en una RAM, sino en la fuerza de sinapsis entre células de Kenyon y neuronas
de salida, modificadas por una regla neuromoduladora dopaminérgica análoga a
la que sustenta el aprendizaje olfativo en la mosca. Demostramos que el
lenguaje es Turing-completo mediante una reducción a máquinas de registros de
Minsky, acotamos la probabilidad de error de su dinámica estocástica y
caracterizamos las órbitas del sistema en el límite determinista. Aportamos
un intérprete en Python, de código abierto y reproducible, que resuelve
programas de ejemplo (impresión de caracteres, «Hola, Mundo» en UTF-8 y suma
de enteros por transferencia sináptica) y cuantifica la actividad evocada en
el resto del cerebro real. Observamos una transición abrupta de fiabilidad en
función de la temperatura neuronal, con fiabilidad perfecta en 40 de 40 ensayos
para $\beta\ge 24$ y fallo total para $\beta\le 12$.

\medskip
\noindent\textbf{Palabras clave:} conectoma, FlyWire, \textit{Drosophila},
lenguajes esotéricos, completitud de Turing, cuerpo en champiñón,
plasticidad dopaminérgica, sistemas dinámicos.
\end{abstract}

\begin{otherlanguage}{english}
\begin{abstract}
We present \dros{}, an esoteric programming language whose execution model
is the brain of a fruit fly rather than a von Neumann machine. Programs run
on the complete adult FlyWire v783 connectome (138{,}639 neurons, 2{,}700{,}513
connections with at least five synapses), with synaptic signs derived from
predicted neurotransmitters. Each instruction is mapped, by an
anatomy-guided recruitment procedure, onto an identified real neuron.
Program memory lives in Kenyon cell--to--mushroom body output neuron
synapses rewritten by a dopaminergic neuromodulatory rule. We prove Turing
completeness by reduction to Minsky register machines, bound the error of
the stochastic dynamics, and release a reproducible Python interpreter.
\end{abstract}
\end{otherlanguage}

%% ============================================================================
\section{Introducción}
\label{sec:intro}

Los lenguajes de programación esotéricos exploran los límites de lo que puede
considerarse cómputo y de cómo puede expresarse~\cite{mateas2005}. Aunque suelen
concebirse como ejercicios lúdicos o estéticos, algunos plantean preguntas
serias sobre el sustrato físico del cálculo. Paralelamente, la neurociencia de
sistemas ha alcanzado un hito: la reconstrucción del diagrama de conexiones
completo del cerebro de una mosca adulta, \Dm{}, por el consorcio
FlyWire~\cite{dorkenwald2024,schlegel2024}. Ese conectoma ya se ha utilizado
para construir modelos computacionales del cerebro entero que predicen
respuestas sensoriomotoras~\cite{shiu2024} y actividad en el sistema
visual~\cite{lappalainen2024}.

Este trabajo une ambas líneas. Nos preguntamos si es posible definir un
lenguaje de programación, con sintaxis, semántica e intérprete, cuyo modelo de
ejecución sea literalmente la propagación de actividad sobre las neuronas
reales identificadas de una mosca, y cuya memoria resida en las sinapsis que
la biología utiliza para aprender. La respuesta que ofrecemos es
\dros{}.

El lenguaje se apoya en tres principios de diseño que conviene explicitar
desde el comienzo. El primero es el \emph{emergentismo}: ninguna neurona
contiene el programa; este solo existe como la trayectoria que recorre la
excitación a lo largo del tiempo. El segundo es la \emph{percepción}: la
ejecución comienza con un estímulo olfativo sobre un glomérulo, y los datos
de entrada se entregan como concentraciones de olor a receptores olfativos
reales. El tercero es la \emph{muerte temporal}: la parada del programa
coincide con el disparo de la Fibra Gigante, la neurona del reflejo de
escape, tras el cual la simulación se detiene.

\paragraph{Contribuciones.} Las aportaciones de este artículo son las siguientes:
\begin{enumerate}[leftmargin=*,itemsep=2pt]
  \item Un modelo dinámico híbrido (\cref{sec:modelo}) que combina neuronas de
        integración y disparo sobre el conectoma FlyWire con un circuito
        implantado de unidades estocásticas de umbral, y una regla de
        plasticidad dopaminérgica en sinapsis célula de Kenyon $\to$ MBON.
  \item Un lenguaje formal, con gramática EBNF y semántica operacional con
        latencias exactas (\cref{sec:lenguaje}).
  \item Una demostración de completitud de Turing por reducción a máquinas
        de Minsky, una cota de error para la dinámica estocástica y una
        tricotomía de órbitas en el límite determinista (\cref{sec:turing}).
  \item Un procedimiento de reclutamiento anatómico que asigna cada papel del
        programa a una neurona real identificada, siguiendo conexiones reales
        cuando existen (\cref{sec:implementacion}).
  \item Un intérprete de código abierto y una evaluación empírica sobre el
        conectoma real (\cref{sec:resultados}).
\end{enumerate}

\paragraph{Alcance.} \dros{} no pretende afirmar que la mosca real ejecute
estos programas. El circuito del programa se implanta sobre neuronas reales,
de forma análoga a una manipulación transgénica, y está protegido por
construcción frente a la actividad natural del resto del cerebro. El valor de
la propuesta reside en su rigor formal, en su fidelidad a la identidad y a la
anatomía de las neuronas empleadas y en su capacidad para visualizar cómo un
cómputo abstracto se propaga por un sistema nervioso real.

%% ============================================================================
\section{Trabajos relacionados}
\label{sec:relacionados}

\paragraph{Cómputo neuronal y completitud de Turing.} McCulloch y Pitts
mostraron que las redes de neuronas de umbral realizan cualquier función
lógica finita~\cite{mcculloch1943}. Siegelmann y Sontag probaron que las redes
recurrentes con pesos racionales y activación saturada lineal son
Turing-completas~\cite{siegelmann1995}. Nuestra construcción se aparta de esta
última: las activaciones son binarias y toda la memoria no acotada reside en
pesos sinápticos enteros, lo que permite una reducción directa a las máquinas
de registros de Minsky~\cite{minsky1967}. La conexión entre sistemas dinámicos
e indecidibilidad fue formalizada, entre otros, por Moore~\cite{moore1990}.

\paragraph{Conectómica y modelos del cerebro entero.} El hemicerebro de
\textit{Drosophila}~\cite{scheffer2020} y, posteriormente, el cerebro completo
FlyWire~\cite{dorkenwald2024}, anotado con más de 8\,000 tipos
celulares~\cite{schlegel2024}, han hecho posible simular el cerebro entero
de la mosca. Shiu \emph{et al.}~\cite{shiu2024} construyeron un modelo de
integración y disparo con fugas sobre FlyWire que reproduce la activación de
la motoneurona MN9 ante estímulos azucarados; su tabla de conectividad es la
que empleamos aquí.

\paragraph{El cuerpo en champiñón como sustrato de memoria.} El cuerpo en
champiñón es el centro del aprendizaje olfativo asociativo en la mosca. Las
células de Kenyon (KC) reciben entradas convergentes y aleatorias de las
neuronas de proyección~\cite{caron2013}; la neurona APL impone una
inhibición global que esparce la codificación~\cite{lin2014}; y las neuronas
dopaminérgicas (DAN) modifican de forma compartimentada las sinapsis entre
KC y neuronas de salida (MBON)~\cite{aso2014,cohn2015,hige2015}. Este es,
precisamente, el mecanismo que \dros{} utiliza como memoria.

\paragraph{Neuronas efectoras identificadas.} El glomérulo DA1 recibe los
receptores Or67d, sensibles a la feromona masculina cVA~\cite{kurtovic2007}.
MN9 controla la extensión de la probóscide~\cite{gordon2009}. DNa02 es una
neurona descendente implicada en el giro durante la marcha, y DNp01 es la
Fibra Gigante que desencadena el escape~\cite{namiki2018,allen2006}.

%% ============================================================================
\section{Modelo dinámico}
\label{sec:modelo}

\subsection{El sustrato: conectoma FlyWire}

Sea $V$ el conjunto de las $N=138\,639$ neuronas de la tabla de conectividad
FlyWire v783 distribuida con el modelo de Shiu \emph{et al.}~\cite{shiu2024}.
Para cada par $(j,i)$ con $c_{ij}\geq 5$ sinapsis de $j$ a $i$ (umbral
estándar de FlyWire), definimos un signo $e_j\in\{+1,-1\}$ según el
neurotransmisor predicho de $j$: la acetilcolina es excitadora y el GABA y el
glutamato, inhibidores. El resultado es un grafo dirigido con 2\,700\,513
aristas. Cada neurona se anota con su tipo celular, superclase, clase y
hemisferio a partir de Schlegel \emph{et al.}~\cite{schlegel2024}. La
\cref{tab:superclases} resume la composición del sustrato.

\begin{table}[t]
\centering
\caption{Composición del sustrato FlyWire v783 por superclase (neuronas
presentes en la tabla de conectividad).}
\label{tab:superclases}
\begin{tabular}{lr@{\hspace{2.5em}}lr}
\toprule
Superclase & Neuronas & Superclase & Neuronas \\
\midrule
óptica                 & 77\,530 & ascendente           & 1\,736 \\
central                & 32\,379 & descendente          & 1\,299 \\
sensorial              & 16\,352 & sensorial ascendente &    581 \\
proyección visual      &  8\,038 & centrífuga visual    &    524 \\
motora                 &    110  & endocrina            &     76 \\
\midrule
\multicolumn{3}{l}{Total (incluidas 14 sin anotar)} & 138\,639 \\
\bottomrule
\end{tabular}
\end{table}

\subsection{Dinámica del cerebro natural}

Los pesos naturales se normalizan por la entrada total de cada neurona:
\begin{equation}
  \tilde w_{ij} \;=\; \frac{e_j\, c_{ij}}{\sum_k c_{ik}},
  \qquad\text{de modo que}\qquad \sum_j \lvert \tilde w_{ij}\rvert = 1 .
  \label{eq:normalizacion}
\end{equation}
Cada neurona natural es una unidad de integración y disparo en tiempo discreto
con potencial $v_i$, fuga $\lambda$, umbral $\theta$ y periodo refractario de
$\rho$ pasos:
\begin{align}
  v_i(t+1) &= \lambda\, v_i(t) + \sum_j \tilde w_{ij}\, s_j(t),
  \label{eq:lif}\\
  s_i(t+1) &= \ind\!\left[\, v_i(t+1)\ge\theta \;\wedge\; r_i(t)=0 \,\right],
  \label{eq:disparo-natural}
\end{align}
donde, tras un disparo, $v_i$ se reinicia a cero y $r_i$ se fija en $\rho$.
La normalización~\eqref{eq:normalizacion} resultó necesaria: con pesos
proporcionales al número bruto de sinapsis, un único disparo de las neuronas
de proyección de DA1 generaba actividad autosostenida en miles de neuronas.
Con $\theta = 0{,}08$, es decir, cuando dispara alrededor del 8\,\% de la
entrada de una neurona, la respuesta evocada es dispersa (unas 350 neuronas)
y se extingue en unos siete pasos.

\subsection{El circuito implantado}

El programa se instala como un circuito de $N_p$ neuronas \emph{reclutadas},
$P\subset V$, que obedecen a unidades estocásticas de umbral. Para $i\in P$,
\begin{align}
  h_i(t) &= \sum_{j} W^{P}_{ij}(t)\, s_j(t) \;+\; b_i \;+\; I_i(t)
            \;+\; \kappa \sum_j \tilde w_{ij}\, s_j(t),
  \label{eq:potencial-implante}\\
  \Pr\!\left[s_i(t+1)=1 \,\middle|\, \mathcal F_t\right] &= \sig\!\big(\beta\, h_i(t)\big),
  \qquad \sig(x) = \frac{1}{1+e^{-x}},
  \label{eq:disparo-implante}
\end{align}
donde $W^P$ son los pesos del implante (enteros), $b_i\in\Z+\tfrac12$ el
sesgo, $I_i(t)$ el estímulo externo, $\beta=1/T$ la inversa de la temperatura
neuronal y $\kappa=\tfrac18$ el coeficiente de \emph{susurro} con el que la
entrada natural influye en la neurona reclutada. Las neuronas reclutadas
conservan sus sinapsis de salida reales, de modo que sus espigas se propagan
al resto del cerebro según~\eqref{eq:lif}.

\subsection{Plasticidad dopaminérgica}

Cada engrama $r\in\mathcal R$ es una sinapsis plástica $K_r\to M_r$ entre una
célula de Kenyon y una MBON, inervada por una dopaminérgica apetitiva
$\PAM_r$ y una aversiva $\PPL_r$. Su peso $w_r(t)\in\N$ evoluciona según
\begin{equation}
  w_r(t+1) = \max\!\Big(0,\; w_r(t) + \eta\,\big[s_{\PAM_r}(t) - s_{\PPL_r}(t)\big]\Big),
  \qquad \eta = 1 .
  \label{eq:plasticidad}
\end{equation}
Equivalentemente, $W^P(t) = W^P_0 + \sum_r w_r(t)\, e_{M_r}e_{K_r}^{\!\top}$.
La regla es neuromoduladora y compartimentada, en consonancia con la
organización del cuerpo en champiñón~\cite{aso2014,cohn2015,hige2015}.

\subsection{Estímulo, lectura y parada}

El estímulo inicial es $I(0)=e_{c_0}$, con $c_0$ la neurona de entrada del
glomérulo soplado. La salida se lee en los instantes $t_k$ en que dispara la
neurona de estrobo DNa02:
\begin{equation}
  y_k = \sum_{j=0}^{7} 2^{j}\, s_{\mathrm{MN}_j}(t_k),
  \label{eq:byte}
\end{equation}
y se interpreta el flujo de bytes como UTF-8. Un segundo canal cuenta las
espigas de MN9 entre estrobos y las imprime en decimal. El tiempo de muerte es
$\tau=\inf\{t\ge1: s_{\GF}(t)=1\}$. La \cref{tab:parametros} recoge los
parámetros empleados.

\begin{table}[t]
\centering
\caption{Parámetros del modelo empleados en todos los experimentos.}
\label{tab:parametros}
\begin{tabular}{llc}
\toprule
Símbolo & Significado & Valor \\
\midrule
$c_{\min}$ & sinapsis mínimas por conexión & 5 \\
$\lambda$  & fuga de membrana natural por paso & 0,5 \\
$\theta$   & umbral natural (fracción de entrada) & 0,08 \\
$\rho$     & periodo refractario natural (pasos) & 2 \\
$\kappa$   & coeficiente de susurro sobre el implante & $1/8$ \\
$\beta$    & inversa de la temperatura del implante & 64 \\
$\eta$     & cuanto de plasticidad dopaminérgica & 1 \\
\bottomrule
\end{tabular}
\end{table}

%% ============================================================================
\section{El lenguaje}
\label{sec:lenguaje}

\subsection{Morfología}

Un programa \dros{} es un mapa de glomérulos. Cada línea define un glomérulo,
es decir, un estado del programa, seguido de una secuencia de acciones y de un
desenlace. Si el nombre del glomérulo coincide con uno real del lóbulo antenal
(de entre los 53 con receptores anotados en FlyWire), el estado se asigna a su
neurona de proyección real. El \cref{lst:ejemplo} muestra la forma general.

\begin{lstlisting}[style=dros,caption={Forma general de un programa.},label={lst:ejemplo}]
%cepa Canton-S            ;; semilla del ruido neuronal
$ memoria = 3             ;; engrama: sinapsis KC -> MBON con fuerza 3
~~~> (DM1)                ;; soplo inicial I(0)
(DM1) +memoria !'a' ~> (VA2)
(VA2) ?memoria ~> (DM1) | (DL5)
(DL5) x_x                 ;; muerte
\end{lstlisting}

\subsection{Gramática}

La \cref{fig:ebnf} define la sintaxis en EBNF. Se admiten alias Unicode
(\S{} para \texttt{\$}, $\approx\approx>$ para \texttt{\textasciitilde\textasciitilde\textasciitilde>},
$\rightsquigarrow$ para \texttt{\textasciitilde>} y \dag{} para \texttt{x\_x}).

\begin{figure}[t]
\begin{lstlisting}[basicstyle=\ttfamily\footnotesize]
programa   = { linea } ;
linea      = [ sentencia ] [ ";;" comentario ] salto ;
sentencia  = directiva | engrama | soplo | glomerulo ;
directiva  = "%" ( "cepa" | "temperatura" ) valor ;
engrama    = ( "§" | "$" ) nombre [ "=" natural ] ;
soplo      = ( "~~~>" | "≈≈>" ) glom ;
glomerulo  = glom { accion } desenlace ;
glom       = "(" nombre ")" ;
accion     = "+" nombre | "-" nombre | "<" nombre | "!" emision ;
emision    = caracter | cadena | natural | "*" ;      (* natural en [0,255] *)
desenlace  = flecha glom
           | "?" nombre flecha glom "|" glom
           | "x_x" | "†" ;
flecha     = "~>" | "⇝" ;
nombre     = letra { letra | digito | "_" } ;        (* distinto de "x_x" *)
\end{lstlisting}
\caption{Gramática de \dros{} en notación EBNF.}
\label{fig:ebnf}
\end{figure}

\paragraph{Semántica estática.} Los nombres de glomérulo son únicos; todo
destino debe estar definido; hay a lo sumo un soplo y, en su ausencia, se
sopla el primer glomérulo. Un engrama usado sin declarar nace con $w=0$. La
anatomía impone una restricción adicional: el conectoma contiene solo 16
neuronas PPL1, por lo que un programa admite como máximo 16 engramas, y cada
prueba \texttt{?r} consume cuatro de las 514 neuronas locales del cuerno
lateral, por lo que admite como máximo 128 pruebas. Ambos límites dejan sitio
para una máquina universal (\cref{sec:turing}).

\subsection{Semántica operacional}

Cada glomérulo se despliega en una cadena de micro-estados, uno por acción,
y cada micro-estado ocupa una neurona reclutada. La \cref{tab:semantica}
especifica, para cada construcción, el circuito implantado, las neuronas
reales que lo ejecutan y su latencia exacta en pasos.

\begin{table}[t]
\centering
\caption{Semántica operacional: construcción, circuito implantado, neuronas
reales reclutadas y latencia. $C$ denota la neurona del micro-estado activo y
$c$ la concentración de entrada.}
\label{tab:semantica}
\small
\begin{tabularx}{\linewidth}{@{}l X l c@{}}
\toprule
Construcción & Circuito & Neuronas reales & Latencia \\
\midrule
\texttt{+r}        & $C\to\PAM_r$; $C\to$ siguiente & PAM & 1 \\
\texttt{-r}        & $C\to\PPL_r$; $C\to$ siguiente & PPL1 & 1 \\
\texttt{<r}        & ráfaga de $c$ espigas en $\ORN_r\to\PAM_r$; retorno & ORN, ALLN & $c+4$ \\
\texttt{!b}        & $C\to\{\mathrm{MN}_j: b_j=1\}\cup\{\mathrm{DNa02}\}$ & motoneuronas, DNa02 & 1 por byte \\
\texttt{!*}        & $C\to\mathrm{MN9}$ & MN9 & 1 \\
\texttt{\textasciitilde> (X)} & fusionado con la última acción (relevo si no hay acciones) & --- & 0 o 1 \\
\texttt{?r \textasciitilde> (A) | (B)} & $K_r$, $M_r$, dos retardos y dos compuertas (\cref{fig:prueba}) & KC, MBON, LHLN & 4 \\
\texttt{x\_x}      & $C\to\GF$ & DNp01 & 1 \\
\bottomrule
\end{tabularx}
\end{table}

\subsection{Entrada y salida}

La entrada es una cola de concentraciones naturales. Cuando dispara el
micro-estado de \texttt{<r}, el entorno entrega $c$ espigas consecutivas a
$\ORN_r$, cada una de las cuales activa $\PAM_r$ e incrementa $w_r$ en una
unidad según~\eqref{eq:plasticidad}; al terminar, una neurona local del
lóbulo antenal devuelve el control al programa. Si la cola se ha agotado,
$c=0$. La salida se decodifica según~\eqref{eq:byte} y mediante el conteo de
MN9.

%% ============================================================================
\section{Completitud de Turing y propiedades dinámicas}
\label{sec:turing}

\subsection{Puertas lógicas en el cuerpo en champiñón}

Con $\Theta$ la función escalón, las puertas elementales se escriben como
neuronas de umbral con sesgo semientero:
\begin{equation}
\begin{aligned}
  \mathrm{Y}(x,y)&=\Theta\!\left(x+y-\tfrac32\right), &\qquad
  \mathrm{O}(x,y)&=\Theta\!\left(x+y-\tfrac12\right),\\
  \mathrm{NO}(x)&=\Theta\!\left(\tfrac12-x\right), &\qquad
  \mathrm{NOY}(x,y)&=\Theta\!\left(\tfrac32-x-y\right).
\end{aligned}
  \label{eq:puertas}
\end{equation}
La puerta Y corresponde a la detección de coincidencias de una célula de
Kenyon; la NO, a la inhibición GABAérgica de la neurona APL; y la NOY, a una
MBON con actividad tónica que recibe inhibición. Como NOY es funcionalmente
completa, todo circuito booleano finito es realizable~\cite{mcculloch1943}.

\subsection{Cota de error}

\begin{lema}[Margen]
\label{lem:margen}
En todo circuito compilado y para toda neurona reclutada $i$, se cumple
$\lvert h_i(t)\rvert \ge \tfrac38$.
\end{lema}
\begin{proof}
Los pesos del implante y el estímulo son enteros, y los sesgos pertenecen a
$\Z+\tfrac12$; por tanto, la parte del implante en~\eqref{eq:potencial-implante}
dista al menos $\tfrac12$ de cero. Por~\eqref{eq:normalizacion},
$\lvert\kappa\sum_j\tilde w_{ij}s_j\rvert \le \kappa=\tfrac18$. Luego
$\lvert h_i\rvert\ge\tfrac12-\tfrac18=\tfrac38$.
\end{proof}

\begin{corolario}
\label{cor:error}
La probabilidad de que, en $\tau$ pasos, la trayectoria estocástica del
implante difiera de su límite determinista ($\beta\to\infty$) satisface
\begin{equation}
  \Pr[\text{desviación}] \;\le\; \tau\, N_p\, \sig\!\left(-\tfrac{3\beta}{8}\right)
  \;\le\; \tau\, N_p\, e^{-3\beta/8}.
\end{equation}
\end{corolario}
\begin{proof}
Por el \cref{lem:margen}, cada neurona reclutada se aparta del valor
determinista con probabilidad a lo sumo $\sig(-3\beta/8)$ en cada paso. La
cota se sigue de la desigualdad de la unión sobre $N_p$ neuronas y $\tau$ pasos.
\end{proof}

Con $\beta=64$, la cota por neurona y paso es $\sig(-24)\approx3{,}8\times10^{-11}$.

\subsection{Reducción a máquinas de Minsky}

Una máquina de Minsky $\mathcal M=(Q,q_0,\delta)$ con registros
$r_1,\dots,r_m\in\N$ ejecuta instrucciones de tres tipos:
$\mathrm{INC}(r,q')$; $\mathrm{JZDEC}(r,q',q'')$, que resta uno y salta a
$q'$ si $r>0$ o salta a $q''$ si $r=0$; y $\mathrm{HALT}$. Con $m=2$ ya son
universales~\cite{minsky1967}.

\begin{definicion}[Traducción]
\label{def:traduccion}
A cada estado $q$ se le asocia un glomérulo con neurona de control $C_q$
(sesgo $-\tfrac12$) y a cada registro $r$ un engrama. $\mathrm{INC}(r,q')$ se
traduce en $C_q\to\PAM_r$ y $C_q\to C_{q'}$. $\mathrm{HALT}$ se traduce en
$C_q\to\GF$. $\mathrm{JZDEC}(r,q',q'')$ se traduce en el circuito de la
\cref{fig:prueba}: $C_q\to K_r$, la cadena de retardo $C_q\to A_1\to A_2$ y
dos compuertas
\begin{equation}
  G_{\neq0}=\Theta\!\left(A_2+M_r-\tfrac32\right),\qquad
  G_{=0}=\Theta\!\left(A_2-M_r-\tfrac12\right),
\end{equation}
con $M_r=\Theta\!\left(w_r K_r-\tfrac12\right)$, $G_{\neq0}\to\{\PPL_r,C_{q'}\}$
y $G_{=0}\to C_{q''}$.
\end{definicion}

\begin{figure}[t]
\centering
\begin{tikzpicture}[
  >={Stealth[length=2.2mm]},
  neu/.style={circle,draw,thick,minimum size=9mm,inner sep=0pt,font=\footnotesize},
  dan/.style={neu,fill=orange!15,minimum size=12mm,font=\scriptsize},
  gate/.style={neu,fill=blue!8},
  ctl/.style={neu,fill=green!12},
  lab/.style={font=\scriptsize,text=black!70},
  node distance=12mm and 16mm]
  \node[ctl] (C) {$C_q$};
  \node[neu,right=of C] (K) {$K_r$};
  \node[neu,right=22mm of K] (M) {$M_r$};
  \node[neu,below=of C] (A1) {$A_1$};
  \node[neu,right=of A1] (A2) {$A_2$};
  \node[gate,right=22mm of A2,yshift=6mm] (Gnz) {$G_{\neq0}$};
  \node[gate,right=22mm of A2,yshift=-10mm] (Gz) {$G_{=0}$};
  \node[ctl,right=18mm of Gnz] (Q1) {$C_{q'}$};
  \node[ctl,right=18mm of Gz] (Q2) {$C_{q''}$};
  \node[dan,above=10mm of K] (PAM) {$\PAM_r$};
  \node[dan,above=10mm of M] (PPL) {$\PPL_r$};
  \draw[->,thick] (C) -- (K);
  \draw[->,very thick,red!70!black] (K) -- node[below,lab]{$w_r(t)$} (M);
  \draw[->,thick] (C) -- (A1);
  \draw[->,thick] (A1) -- (A2);
  \draw[->,thick] (A2) -- (Gnz);
  \draw[->,thick] (A2) -- (Gz);
  \draw[->,thick] (M) -- node[right,lab]{$+1$} (Gnz);
  \draw[-{Bar[width=2.5mm]},thick] (M) to[bend left=35] node[right,lab,pos=0.7]{$-1$} (Gz);
  \draw[->,thick] (Gnz) -- (Q1);
  \draw[->,thick] (Gz) -- (Q2);
  \draw[->,thick,dashed] (Gnz.east) to[out=15,in=0,looseness=1.6] (PPL.east);
  \draw[->,dashed,orange!70!black] (PAM.south east) to[bend left=10] ($ (K)!0.5!(M) + (0,0.8mm) $);
  \draw[->,dashed,orange!70!black] (PPL.south) -- ($ (K)!0.62!(M) + (0,0.8mm) $);
  \node[lab,align=center,below=1mm of A2] {LHLN};
  \node[lab,align=center,below=1mm of A1] {LHLN};
  \node[lab,above=0.5mm of C] {PN / KC};
\end{tikzpicture}
\caption{Circuito implantado para la instrucción de prueba
\texttt{?r \textasciitilde> (A) | (B)}. La célula de Kenyon $K_r$ lee la
sinapsis plástica $w_r$ (rojo); la MBON $M_r$ responde solo si $w_r\ge1$. Las
compuertas del cuerno lateral combinan esa respuesta con la señal retardada
del estado. Las dopaminérgicas (naranja) modulan la sinapsis del engrama.}
\label{fig:prueba}
\end{figure}

\begin{lema}[Unicidad del foco]
\label{lem:foco}
En el límite determinista, si en el instante de inicio de una instrucción
exactamente una neurona de control está activa y todas las neuronas
transitorias están en reposo, la misma propiedad se cumple al inicio de la
instrucción siguiente.
\end{lema}
\begin{proof}
Por inspección de cada traducción: toda neurona de control proyecta
exclusivamente a su sucesora o a su circuito de prueba, y en este último solo
una de las compuertas $G_{\neq0}$, $G_{=0}$ puede superar su umbral, puesto
que requieren valores opuestos de $M_r$. Las neuronas transitorias carecen de
autoconexiones y vuelven al reposo tras un paso.
\end{proof}

\begin{lema}[Coherencia temporal]
\label{lem:coherencia}
Toda lectura de $w_r$ efectuada por una instrucción observa el efecto de
todas las escrituras de las instrucciones anteriores.
\end{lema}
\begin{proof}
Un $\mathrm{INC}$ que comienza en $t$ activa $\PAM_r$ en $t+1$ y, por
\eqref{eq:plasticidad}, modifica el peso efectivo en $t+2$. Una prueba que
comienza en $t'\ge t+1$ activa $K_r$ en $t'+1\ge t+2$, que lee el peso ya
actualizado. Una prueba iniciada en $t'$ decrementa mediante $G_{\neq0}(t'+3)
\to\PPL_r(t'+4)$, con efecto en $t'+5$, mientras que la siguiente instrucción
comienza en $t'+4$ y lee, a lo sumo, en $t'+5$.
\end{proof}

\begin{lema}[Capacidad]
\label{lem:capacidad}
La traducción de una máquina de Minsky con $m$ registros y $b$ instrucciones
$\mathrm{JZDEC}$ puede implantarse en FlyWire v783 si $m\le16$ y $b\le128$.
\end{lema}
\begin{proof}
Cada registro consume una célula de Kenyon, una MBON, una PAM, una PPL1 y un
receptor olfativo; cada $\mathrm{JZDEC}$, cuatro neuronas locales del cuerno
lateral (dos retardos y dos compuertas); cada estado, una neurona de
proyección o una célula de Kenyon; los efectores, once neuronas fijas. Las
anotaciones de FlyWire contienen 5\,177 células de Kenyon, 96 MBON, 307 PAM,
16 PPL1 y 514 neuronas locales del cuerno lateral. Las restricciones que
primero se agotan son $m\le16$ (PPL1) y $4b\le514$ (cuerno lateral); las
demás quedan muy lejos de agotarse con $m\le16$ y $b\le128$.
\end{proof}

\begin{teorema}[Simulación de máquinas de Minsky]
\label{thm:simulacion}
Para toda máquina de Minsky $\mathcal M$ que cumpla las cotas del
\cref{lem:capacidad} existe un programa $P_{\mathcal M}$ de \dros{} tal que,
en el límite $\beta\to\infty$, $P_{\mathcal M}$ provoca el disparo de la
Fibra Gigante si y solo si $\mathcal M$ se detiene, y en cada macro-paso los
pesos $w_r$ coinciden con los registros de $\mathcal M$.
\end{teorema}
\begin{proof}
Por inducción sobre el número de macro-pasos. La configuración $(q,\vec n)$
de $\mathcal M$ se corresponde con el estado en que $C_q$ es la única neurona
de control activa y $\vec w=\vec n$. El caso base lo proporciona el soplo
$I(0)=e_{C_{q_0}}$ con pesos iniciales declarados. El paso inductivo se sigue
de la \cref{def:traduccion} y de los \cref{lem:foco,lem:coherencia}. El
\cref{lem:capacidad} garantiza que el implante existe.
\end{proof}

\begin{corolario}[Completitud de Turing]
\label{thm:turing}
\dros{} es Turing-completo.
\end{corolario}
\begin{proof}
Un cerebro finito no puede alojar la traducción de \emph{cualquier} máquina de
Minsky, pero no lo necesita: basta con alojar una máquina universal fija.
Korec construyó una máquina de registros universal, $U_{32}$, con 32
instrucciones y 8 registros, cuyas instrucciones son incrementos y
decrementos con salto condicionado al cero, es decir, $\mathrm{INC}$ y
$\mathrm{JZDEC}$~\cite{korec1996}. Tiene como mucho 32 instrucciones
$\mathrm{JZDEC}$ y 8 registros, dentro de las cotas del \cref{lem:capacidad};
por el \cref{thm:simulacion}, \dros{} la simula, y con ella a cualquier
máquina de Turing codificada en sus registros de entrada. La memoria no
acotada reside en los pesos $w_r\in\N$, que el modelo no limita.
\end{proof}

\paragraph{Verificación empírica.} El guion \texttt{minsky.py} genera
máquinas de Minsky aleatorias (de 2 a 12 estados, de 1 a 3 registros), las
traduce según la \cref{def:traduccion}, las ejecuta sobre FlyWire v783 con
$\beta=64$ y compara los registros finales con una simulación directa. En
6\,000 máquinas generadas con dos semillas, las 2\,707 que se detienen en
menos de 2\,000 pasos coincidieron todas con la referencia. Un programa con
128 pruebas se implanta y termina; con 129, el reclutamiento falla por falta
de neuronas del cuerno lateral, como predice el \cref{lem:capacidad}.

\begin{observacion}
Para $\beta$ finito, $P_{\mathcal M}$ es una máquina probabilística cuya
probabilidad de error en ejecuciones de $\tau$ pasos está acotada por el
\cref{cor:error}. La actividad natural del resto del cerebro no altera este
resultado gracias al \cref{lem:margen}.
\end{observacion}

\subsection{Órbitas en el límite determinista}

\begin{proposicion}[Tricotomía orbital]
\label{prop:tricotomia}
Con entrada finita y $\beta=\infty$, toda órbita del implante cumple
exactamente una de las condiciones siguientes: (i) alcanza el disparo de la
Fibra Gigante; (ii) es eventualmente periódica; o (iii) tiene
$\sup_t\max_r w_r(t)=\infty$.
\end{proposicion}
\begin{proof}
El estado del implante es $(s_P,\vec w,\text{agenda})$. Una vez consumida la
entrada, la agenda sensorial está acotada. Si los pesos permanecen acotados,
la órbita recorre un conjunto finito de estados y, por el principio del
palomar, es eventualmente periódica.
\end{proof}

Los casos (i) y (ii) son detectables por observación, y el intérprete los
reconoce como \emph{muerte}, \emph{coma} (ciclo del estado nulo) o
\emph{ciclo}. El caso (iii) es el único donde reside la indecidibilidad del
problema de la parada~\cite{moore1990}: la memoria no acotada, y con ella lo
indecidible, reside exclusivamente en las sinapsis.

%% ============================================================================
\section{Implementación}
\label{sec:implementacion}

\subsection{Arquitectura}

El intérprete \texttt{drosophila.py}, de unas 1\,100 líneas y dependiente de
\texttt{numpy} (con \texttt{pandas} y \texttt{pyarrow} solo para la primera
carga), se organiza en cinco etapas (\cref{fig:arquitectura}).

\begin{figure}[t]
\centering
\begin{tikzpicture}[
  >={Stealth[length=2mm]},
  caja/.style={draw,thick,rounded corners=2pt,align=center,
               minimum height=11mm,text width=22mm,font=\footnotesize,fill=black!3},
  dato/.style={draw,dashed,rounded corners=2pt,align=center,
               text width=22mm,font=\scriptsize,fill=yellow!8},
  node distance=6mm]
  \node[caja] (lex) {Lexer y\\parser};
  \node[caja,right=of lex] (rec) {Reclutamiento\\anatómico};
  \node[caja,right=of rec] (imp) {Implante\\$W^P$, $I(0)$};
  \node[caja,right=of imp] (din) {Dinámica\\$t\to t+1$};
  \node[caja,right=of din] (dec) {Lectura motora\\y muerte};
  \node[dato,above=7mm of lex] (src) {código\\\texttt{.dros}};
  \node[dato,above=7mm of rec] (fw) {FlyWire v783\\conectividad\\y anotaciones};
  \node[dato,above=7mm of din] (olor) {olor de entrada\\(\texttt{-{}-olor})};
  \node[dato,above=7mm of dec] (out) {bytes UTF-8\\y epitafio};
  \draw[->,thick] (lex) -- (rec);
  \draw[->,thick] (rec) -- (imp);
  \draw[->,thick] (imp) -- (din);
  \draw[->,thick] (din) -- (dec);
  \draw[->] (src) -- (lex);
  \draw[->] (fw) -- (rec);
  \draw[->] (fw.east) to[bend left=12] (din.north west);
  \draw[->] (olor) -- (din);
  \draw[->] (dec) -- (out);
\end{tikzpicture}
\caption{Arquitectura del intérprete. El conectoma real interviene tanto en
el reclutamiento de neuronas como en la dinámica de cada paso.}
\label{fig:arquitectura}
\end{figure}

\begin{enumerate}[leftmargin=*,itemsep=2pt]
  \item \textbf{Análisis léxico y sintáctico.} Tokenizador por patrones
        ordenados y analizador descendente línea a línea, con diagnósticos
        que indican línea y columna.
  \item \textbf{Carga del conectoma.} Descarga única (unos 135\,MB) de la
        conectividad de Shiu \emph{et al.}~\cite{shiu2024} y de las
        anotaciones de Schlegel \emph{et al.}~\cite{schlegel2024}; se
        almacena en caché en formato \texttt{npz} comprimido (11\,MB).
  \item \textbf{Reclutamiento} (\cref{alg:reclutamiento}).
  \item \textbf{Dinámica}: integración de
        \eqref{eq:lif}--\eqref{eq:plasticidad} en cada paso.
  \item \textbf{Decodificación y parada}, con detección de muerte, coma y
        ciclos mediante una tabla de estados visitados.
\end{enumerate}

\subsection{Reclutamiento anatómico}

Cada papel del programa se asigna a una neurona libre del tipo
correspondiente. Cuando existen conexiones reales pertinentes, se elige la
candidata con más sinapsis reales hacia (o desde) las neuronas ya
reclutadas; en caso contrario, se toma la primera por hemisferio (derecho
primero) e identificador. El \cref{alg:reclutamiento} detalla el caso del
engrama. Las neuronas de control se asignan del mismo modo: la primera de
cada glomérulo con nombre real es su neurona de proyección, y las siguientes
son las células de Kenyon que más sinapsis reales reciben de ella.

\begin{algorithm}[t]
\caption{Reclutamiento anatómico de un engrama $r$.}
\label{alg:reclutamiento}
\begin{algorithmic}[1]
\Require conectoma $(V,c)$; conjunto de neuronas libres $L\subseteq V$
\Ensure neuronas $(K_r,M_r,\PAM_r,\PPL_r,\ORN_r)$
\State $E\gets\{(k,m): k\in\KC,\ m\in\MBON,\ c_{mk}\ge c_{\min}\}$ ordenado por $c_{mk}$ decreciente
\State $(K_r,M_r)\gets$ primer par de $E$ con $k,m\in L$
\State $\PAM_r\gets\arg\max_{p\in\PAM\cap L} c_{M_r p}$ \Comment{inervación real de $M_r$}
\State $\PPL_r\gets\arg\max_{q\in\PPL\cap L} c_{M_r q}$
\State $\ORN_r\gets$ primer receptor libre del glomérulo asociado a $r$
\State $L\gets L\setminus\{K_r,M_r,\PAM_r,\PPL_r,\ORN_r\}$
\State \Return $(K_r,M_r,\PAM_r,\PPL_r,\ORN_r)$
\end{algorithmic}
\end{algorithm}

\subsection{Eficiencia}

Las aristas se almacenan en formato disperso ordenado por neurona
presináptica. En cada paso solo se propagan las espigas de las neuronas
activas, recogiendo sus segmentos de aristas mediante índices vectorizados y
acumulando con \texttt{numpy.bincount}. El coste por paso es, por tanto,
proporcional al número de sinapsis salientes de las neuronas activas y no al
tamaño total del conectoma. El \cref{lst:paso} reproduce el núcleo del bucle.

\begin{lstlisting}[style=py,caption={Núcleo de un paso de simulación.},label={lst:paso}]
activas = np.flatnonzero(S)
entrada_natural = C.empujar(activas)          # espigas por sinapsis reales
h = (np.bincount(P.post_local, weights=W * S[P.pre], minlength=len(rec))
     + P.sesgo + I + SUSURRO * entrada_natural[rec])
S_rec = self._disparar(h, rng)                # implante: Bernoulli(sigma(beta*h))
v = FUGA_NATURAL * v + entrada_natural        # cerebro natural: integrar y disparar
bloqueadas = refr > 0
v[bloqueadas] = 0.0; refr[bloqueadas] -= 1
S_sig = (v >= UMBRAL_NATURAL) & natural
v[S_sig] = 0.0; refr[S_sig] = REFRACTARIO
delta = S[P.dan_up].astype(float) - S[P.dan_down].astype(float)
W[P.arista_engrama] = np.maximum(0.0, W[P.arista_engrama] + eta * delta)
\end{lstlisting}

%% ============================================================================
\section{Resultados}
\label{sec:resultados}

Todos los experimentos se realizaron sobre FlyWire v783 con los parámetros de
la \cref{tab:parametros}, en un MacBook Pro con Apple M2 Pro. La primera carga
del conectoma (descarga y construcción) requirió unos 10\,s; las cargas
posteriores desde caché, 0,3\,s.

\subsection{Programas de ejemplo}

La \cref{tab:resultados} resume la ejecución de los programas del
\cref{app:programas}. Todos producen la salida esperada. La columna
\emph{despertadas} recoge el número de neuronas naturales, no reclutadas, que
dispararon al menos una vez durante la ejecución; la \emph{fidelidad}, la
proporción de sinapsis del implante que ya existen en el conectoma real.

\begin{table}[t]
\centering
\caption{Ejecución de programas de ejemplo sobre FlyWire v783. $\tau$: paso de
muerte; $N_p$: neuronas reclutadas.}
\label{tab:resultados}
\small
\begin{tabular}{@{}llllrrrr@{}}
\toprule
Programa & Entrada & Salida & $\tau$ & $N_p$ & Despertadas & Fidelidad & Tiempo \\
\midrule
\texttt{efe.dros}  & ---      & \texttt{F}            &    3 & 13 &     76 & 1/6   & <0,1\,s \\
\texttt{hola.dros} & ---      & \texttt{¡Hola, Mundo!} &   17 & 27 &    860 & 1/83  & <0,1\,s \\
\texttt{suma.dros} & 0, 0     & \texttt{suma: 0}      &   26 & 51 & 3\,248 & 8/84  & <0,1\,s \\
\texttt{suma.dros} & 0, 5     & \texttt{suma: 5}      &   80 & 51 & 3\,661 & 8/84  & <0,1\,s \\
\texttt{suma.dros} & 3, 4     & \texttt{suma: 7}      &   87 & 51 & 3\,658 & 8/84  & <0,1\,s \\
\texttt{suma.dros} & 12, 30   & \texttt{suma: 42}     &  427 & 51 & 3\,970 & 8/84  & 0,21\,s \\
\texttt{suma.dros} & 200, 300 & \texttt{suma: 500}    & 4\,525 & 51 & 3\,971 & 8/84 & 2,2\,s \\
\texttt{longitud.dros} & \texttt{mosca} & \texttt{bytes: 5} & 2\,747 & 55 & 2\,173 & 10/98 & 1,2\,s \\
\bottomrule
\end{tabular}
\end{table}

\subsection{Traza de la primera palabra}

El programa \texttt{efe.dros} sopla sobre DA1 y emite el carácter \texttt{F}
($70=2^6+2^2+2^1$). El reparto real asigna el estado inicial a la neurona de
proyección DA1\_lPN (id 720575940604407468), la instrucción de muerte a una
célula de Kenyon $\gamma$ (KCg-m), el estrobo a DNa02 y la muerte a DNp01. En
$t=1$ dispara DA1\_lPN; en $t=2$ disparan las motoneuronas de los bits 1, 2 y
6 junto con DNa02, se lee \texttt{F} y 64 neuronas naturales responden a la
activación del glomérulo feromonal, 48 de ellas células de Kenyon; en $t=3$
dispara la Fibra Gigante. La traza completa figura en el \cref{app:traza}.

\subsection{Suma por transferencia sináptica}

El programa \texttt{suma.dros} lee dos concentraciones en los engramas
\texttt{alfa} y \texttt{beta}, transfiere \texttt{beta} a \texttt{alfa} unidad a
unidad y vacía \texttt{alfa} a través de MN9. El reclutamiento seleccionó, a
partir de la conectividad real, la pareja KCg-s2~$\to$~MBON05 como sinapsis
de ambos engramas, la dopaminérgica PAM08 como escritora y PPL101 y PPL108
como borradoras (\cref{tab:reparto}). En la traza con entrada $(2,1)$ se
observa cómo la MBON05 responde a la célula de Kenyon cuando el engrama vale
1 (paso 14) y permanece en silencio cuando vale 0 (paso 19), lo que dirige el
flujo de control a través de las compuertas del cuerno lateral.

\begin{table}[t]
\centering
\caption{Reparto parcial de \texttt{suma.dros}: papel del programa y neurona
real reclutada (tipo celular FlyWire, hemisferio e identificador).}
\label{tab:reparto}
\small
\begin{tabular}{@{}lllr@{}}
\toprule
Papel & Tipo celular & Hemisferio & Identificador FlyWire \\
\midrule
KC$\cdot$alfa         & KCg-s2    & izq. & 720575940620684014 \\
MBON$\cdot$alfa       & MBON05    & der. & 720575940630496374 \\
PAM$\cdot$alfa        & PAM08     & izq. & 720575940621617286 \\
PPL1$\cdot$alfa       & PPL108    & der. & 720575940615750562 \\
ORN$\cdot$alfa        & ORN\_DM2  & der. & 720575940603356128 \\
PPL1$\cdot$beta       & PPL101    & der. & 720575940617691170 \\
estado \texttt{(DM1)} & DM1\_lPN  & der. & 720575940619071005 \\
estado \texttt{(VA2)} & VA2\_adPN & der. & 720575940609460491 \\
estrobo               & DNa02     & der. & 720575940604737708 \\
probóscide            & MN9       & der. & 720575940660219265 \\
muerte                & DNp01     & der. & 720575940632499757 \\
\bottomrule
\end{tabular}
\end{table}

La \cref{fig:actividad} muestra el número de neuronas naturales activas en
cada paso para la entrada $(3,4)$. Pueden distinguirse las cinco fases del
programa: percepción, transferencia, escritura, vaciado y muerte. Los picos
de actividad (hasta 637 neuronas) coinciden con la activación de receptores
olfativos y de neuronas de proyección, que reclutan de forma dispersa
células de Kenyon del cáliz.

\begin{figure}[t]
\centering
\begin{tikzpicture}
\begin{axis}[
  width=0.88\linewidth, height=5.6cm, scale only axis=false,
  xlabel={paso $t$}, ylabel={neuronas naturales activas},
  xmin=0, xmax=88, ymin=0, ymax=700,
  grid=major, grid style={black!10},
  tick label style={font=\footnotesize}, label style={font=\small},
  /pgf/number format/use comma,
  /pgf/number format/1000 sep={\,},
]
  \fill[green!8]  (axis cs:0.5,0)  rectangle (axis cs:15.5,700);
  \fill[blue!6]   (axis cs:15.5,0) rectangle (axis cs:39.5,700);
  \fill[orange!8] (axis cs:39.5,0) rectangle (axis cs:46.5,700);
  \fill[red!5]    (axis cs:46.5,0) rectangle (axis cs:84.5,700);
  \fill[black!8]  (axis cs:84.5,0) rectangle (axis cs:87.5,700);
  \node[font=\scriptsize] at (axis cs:8,670)  {percepción};
  \node[font=\scriptsize] at (axis cs:27.5,670) {transferencia};
  \node[font=\scriptsize,rotate=90] at (axis cs:42.5,560) {escritura};
  \node[font=\scriptsize] at (axis cs:65,670) {vaciado por MN9};
  \node[font=\scriptsize,rotate=90] at (axis cs:86,580) {muerte};
  \addplot[thick,blue!60!black,mark=*,mark size=0.9pt] coordinates {
  (1,0) (2,462) (3,101) (4,147) (5,78) (6,23) (7,4) (8,54) (9,22) (10,637)
  (11,138) (12,234) (13,175) (14,127) (15,65) (16,73) (17,505) (18,224)
  (19,222) (20,208) (21,277) (22,521) (23,271) (24,300) (25,217) (26,227)
  (27,492) (28,270) (29,259) (30,214) (31,243) (32,595) (33,220) (34,247)
  (35,240) (36,290) (37,507) (38,188) (39,208) (40,261) (41,321) (42,165)
  (43,152) (44,145) (45,98) (46,71) (47,90) (48,201) (49,133) (50,101)
  (51,131) (52,196) (53,120) (54,76) (55,95) (56,206) (57,195) (58,119) (59,124) (60,131) (61,112) (62,137) (63,107) (64,200) (65,133) (66,112) (67,187) (68,149) (69,105) (70,73) (71,100) (72,260) (73,152) (74,105) (75,120) (76,144) (77,168) (78,89) (79,94) (80,199) (81,140) (82,166) (83,138) (84,135) (85,102) (86,141) (87,97)};
\end{axis}
\end{tikzpicture}
\caption{Actividad evocada en el cerebro real FlyWire durante la ejecución
de \texttt{suma.dros} con entrada $(3,4)$. Las bandas indican las fases del
programa. En total despertaron 3\,658 neuronas distintas.}
\label{fig:actividad}
\end{figure}

\subsection{Puertas lógicas}

La \cref{tab:puertas} recoge la verificación empírica de las
puertas~\eqref{eq:puertas} con neuronas estocásticas a $\beta=64$: 10\,000
ensayos por fila, 200\,000 en total, sin ningún error.

\begin{table}[t]
\centering
\caption{Tablas de verdad empíricas de las puertas de umbral estocásticas
($\beta=64$, $10^4$ ensayos por combinación de entradas).}
\label{tab:puertas}
\small
\begin{tabular}{@{}llccccr@{}}
\toprule
Puerta & Correlato anatómico & 00 & 01 & 10 & 11 & Errores \\
\midrule
Y    & célula de Kenyon (coincidencia)  & 0 & 0 & 0 & 1 & 0 \\
O    & MBON excitada                    & 0 & 1 & 1 & 1 & 0 \\
NO   & inhibición GABAérgica (APL)      & 1 & 1 & 0 & 0 & 0 \\
NOY  & MBON tónica inhibida             & 1 & 1 & 1 & 0 & 0 \\
XOR  & dos capas (O $\wedge$ NOY)       & 0 & 1 & 1 & 0 & 0 \\
\bottomrule
\end{tabular}
\end{table}

\subsection{Robustez frente a la temperatura}

Ejecutamos \texttt{suma.dros} con entrada $(3,4)$ para distintos valores de
$\beta$ con 40 semillas independientes cada uno. La \cref{tab:temperatura}
revela una transición abrupta: la fiabilidad es perfecta para $\beta\ge24$,
cae al 50\,\% para $\beta=16$ y es nula para $\beta\le12$. En el régimen de
alta temperatura, los fallos típicos son disparos espurios de la Fibra
Gigante (muerte prematura) y emisión de bytes aleatorios. Los resultados son
consistentes con el \cref{cor:error}. Para $\beta=16$, la cota por neurona y
paso es $\sig(-6)\approx2{,}5\times10^{-3}$; con $N_p=51$ y $\tau=87$, el
número esperado de desviaciones por ejecución es de aproximadamente $11$,
por lo que la cota deja de ser informativa. Para $\beta=24$, la cota de la
unión es de $0{,}55$, mientras que no se observó ningún fallo en 40
ejecuciones: la cota es conservadora, porque no toda desviación individual
altera la salida.

\begin{table}[t]
\centering
\caption{Fiabilidad de \texttt{suma.dros} (entrada $(3,4)$) en función de la
inversa de la temperatura del implante. Cota: $\sig(-3\beta/8)$ por neurona y
paso.}
\label{tab:temperatura}
\small
\begin{tabular}{@{}rrrr@{}}
\toprule
$\beta$ & $T$ & Cota por neurona y paso & Ejecuciones correctas \\
\midrule
64 & 0,016 & $3{,}8\times10^{-11}$ & 40/40 \\
32 & 0,031 & $6{,}1\times10^{-6}$  & 40/40 \\
24 & 0,042 & $1{,}2\times10^{-4}$  & 40/40 \\
16 & 0,063 & $2{,}5\times10^{-3}$  & 20/40 \\
12 & 0,083 & $1{,}1\times10^{-2}$  & 0/40 \\
 8 & 0,125 & $4{,}7\times10^{-2}$  & 0/40 \\
\bottomrule
\end{tabular}
\end{table}

%% ============================================================================
\section{Discusión}
\label{sec:discusion}

\paragraph{Fidelidad anatómica.} Las identidades, tipos celulares y
conexiones de salida de todas las neuronas son reales, y el reclutamiento
favorece las conexiones existentes. Aun así, solo 8 de las 84 sinapsis (10\,\%)
del implante de \texttt{suma.dros} existen en el conectoma: el resto actúa
como un transgén. Esta proporción cuantifica el coste de programar un
cerebro que no evolucionó para sumar. Resulta llamativo, sin embargo, que el
reclutamiento guiado por la conectividad asocie de forma espontánea MBON05
con PAM08, dos tipos que comparten compartimento en el lóbulo $\gamma$.

\paragraph{Protección del cómputo.} El coeficiente de susurro $\kappa=\tfrac18$
y la normalización~\eqref{eq:normalizacion} garantizan que la actividad
natural no puede corromper el implante (\cref{lem:margen}). Se trata de una
elección de diseño deliberada: sin ella, el cómputo dependería de la
dinámica natural, que no está calibrada biofísicamente. La contrapartida es
que el cerebro natural observa y reacciona al programa, pero no participa en
sus decisiones.

\paragraph{Limitaciones del modelo natural.} La dinámica natural es una
simplificación en tiempo discreto de un modelo de integración y disparo;
no pretende reproducir tasas de disparo fisiológicas ni el modelo de Shiu
\emph{et al.}~\cite{shiu2024}. En particular, la normalización por entrada
total altera la eficacia relativa de las sinapsis, y en nuestras pruebas no
reprodujo la activación de MN9 por neuronas gustativas de azúcar descrita
por dichos autores.

\paragraph{Restricciones impuestas por la anatomía.} El número de engramas
queda limitado a 16 por el número de neuronas PPL1, y el de pruebas a 128 por
el de neuronas locales del cuerno lateral. Estas restricciones no comprometen
la universalidad, porque la máquina universal de Korec cabe en ellas
(\cref{thm:turing}), pero ilustran cómo la anatomía real se convierte en una
restricción del lenguaje. Una restricción de otro tipo es que la salida solo
puede emitir bytes fijados en el código o el recuento decimal de espigas de
MN9: un byte calculado no puede emitirse directamente, lo que limita la
manipulación de texto a programas como \texttt{longitud.dros}, que leen texto
y producen números.

\paragraph{Trabajo futuro.} Entre las extensiones naturales se cuentan:
permitir que la dinámica natural participe en el cómputo mediante
instrucciones que consulten neuronas no reclutadas; sustituir la dinámica
discreta por un modelo biofísico calibrado; explotar la conectividad real
para compilar programas que maximicen la fidelidad anatómica; y portar el
lenguaje a otros conectomas, como el de la larva de \textit{Drosophila}.

%% ============================================================================
\section{Conclusiones}
\label{sec:conclusiones}

Hemos presentado \dros{}, un lenguaje esotérico Turing-completo que se
ejecuta sobre el conectoma completo de una mosca adulta. Su memoria reside
en sinapsis célula de Kenyon $\to$ MBON reescritas por dopamina, su control
de flujo en neuronas de proyección, células de Kenyon y neuronas del cuerno
lateral, y su salida en motoneuronas identificadas, hasta la Fibra Gigante
que pone fin a la ejecución. Hemos demostrado su completitud de Turing,
acotado su error estocástico y caracterizado sus órbitas, y hemos mostrado
empíricamente que el intérprete ejecuta correctamente programas no triviales
mientras cuantifica la respuesta evocada en miles de neuronas reales.
Además de su interés como objeto lúdico, \dros{} ofrece un banco de pruebas
formal y reproducible para reflexionar sobre la relación entre la estructura
de un conectoma y el cómputo que puede albergar.

%% ============================================================================
\section*{Disponibilidad de código y datos}
El intérprete \texttt{drosophila.py} y los programas de ejemplo están
disponibles en \url{https://github.com/USUARIO/drosophila-lang}. Los datos de
conectividad proceden del repositorio público de Shiu \emph{et al.}
(\url{https://github.com/philshiu/Drosophila_brain_model}) y las anotaciones,
de \url{https://github.com/flyconnectome/flywire_annotations}; su uso está
sujeto a los términos del consorcio FlyWire.

\section*{Agradecimientos}
[Completar: financiación, colaboraciones y, en su caso, declaración sobre el
uso de herramientas de asistencia en la redacción o el desarrollo de
software, conforme a la política de arXiv.]

%% ============================================================================
\begin{thebibliography}{99}
\small

\bibitem{allen2006}
M.~J. Allen, T.~A. Godenschwege, M.~A. Tanouye y P.~Phelan.
\newblock Making an escape: development and function of the \textit{Drosophila}
  giant fibre system.
\newblock \emph{Seminars in Cell \& Developmental Biology}, 17(1):31--41, 2006.

\bibitem{aso2014}
Y.~Aso, D.~Hattori, Y.~Yu, R.~M. Johnston, N.~A. Iyer, T.-T.~B. Ngo, H.~Dionne,
  L.~F. Abbott, R.~Axel, H.~Tanimoto y G.~M. Rubin.
\newblock The neuronal architecture of the mushroom body provides a logic for
  associative learning.
\newblock \emph{eLife}, 3:e04577, 2014.

\bibitem{caron2013}
S.~J.~C. Caron, V.~Ruta, L.~F. Abbott y R.~Axel.
\newblock Random convergence of olfactory inputs in the \textit{Drosophila}
  mushroom body.
\newblock \emph{Nature}, 497:113--117, 2013.

\bibitem{cohn2015}
R.~Cohn, I.~Morantte y V.~Ruta.
\newblock Coordinated and compartmentalized neuromodulation shapes sensory
  processing in \textit{Drosophila}.
\newblock \emph{Cell}, 163(7):1742--1755, 2015.

\bibitem{dorkenwald2024}
S.~Dorkenwald \emph{et al.}
\newblock Neuronal wiring diagram of an adult brain.
\newblock \emph{Nature}, 634:124--138, 2024.

\bibitem{gordon2009}
M.~D. Gordon y K.~Scott.
\newblock Motor control in a \textit{Drosophila} taste circuit.
\newblock \emph{Neuron}, 61(3):373--384, 2009.

\bibitem{hige2015}
T.~Hige, Y.~Aso, M.~N. Modi, G.~M. Rubin y G.~C. Turner.
\newblock Heterosynaptic plasticity underlies aversive olfactory learning in
  \textit{Drosophila}.
\newblock \emph{Neuron}, 88(5):985--998, 2015.

\bibitem{kurtovic2007}
A.~Kurtovic, A.~Widmer y B.~J. Dickson.
\newblock A single class of olfactory neurons mediates behavioural responses to
  a \textit{Drosophila} sex pheromone.
\newblock \emph{Nature}, 446:542--546, 2007.

\bibitem{lappalainen2024}
J.~K. Lappalainen \emph{et al.}
\newblock Connectome-constrained networks predict neural activity across the
  fly visual system.
\newblock \emph{Nature}, 634:1132--1140, 2024.

\bibitem{lin2014}
A.~C. Lin, A.~M. Bygrave, A.~de~Calignon, T.~Lee y G.~Miesenb\"ock.
\newblock Sparse, decorrelated odor coding in the mushroom body enhanced by
  lateral inhibition.
\newblock \emph{Nature Neuroscience}, 17:559--568, 2014.

\bibitem{mateas2005}
M.~Mateas y N.~Montfort.
\newblock A box, darkly: obfuscation, weird languages, and code aesthetics.
\newblock En \emph{Proceedings of the 6th Digital Arts and Culture Conference},
  pp.~144--153, Copenhague, 2005.

\bibitem{mcculloch1943}
W.~S. McCulloch y W.~Pitts.
\newblock A logical calculus of the ideas immanent in nervous activity.
\newblock \emph{Bulletin of Mathematical Biophysics}, 5:115--133, 1943.

\bibitem{minsky1967}
M.~L. Minsky.
\newblock \emph{Computation: Finite and Infinite Machines}.
\newblock Prentice-Hall, Englewood Cliffs, NJ, 1967.

\bibitem{korec1996}
I.~Korec.
\newblock Small universal register machines.
\newblock \emph{Theoretical Computer Science}, 168:267--301, 1996.

\bibitem{moore1990}
C.~Moore.
\newblock Unpredictability and undecidability in dynamical systems.
\newblock \emph{Physical Review Letters}, 64(20):2354--2357, 1990.

\bibitem{namiki2018}
S.~Namiki, M.~H. Dickinson, A.~M. Wong, W.~Korff y G.~M. Card.
\newblock The functional organization of descending sensory-motor pathways in
  \textit{Drosophila}.
\newblock \emph{eLife}, 7:e34272, 2018.

\bibitem{scheffer2020}
L.~K. Scheffer \emph{et al.}
\newblock A connectome and analysis of the adult \textit{Drosophila} central
  brain.
\newblock \emph{eLife}, 9:e57443, 2020.

\bibitem{schlegel2024}
P.~Schlegel \emph{et al.}
\newblock Whole-brain annotation and multi-connectome cell typing of
  \textit{Drosophila}.
\newblock \emph{Nature}, 634:139--152, 2024.

\bibitem{shiu2024}
P.~K. Shiu \emph{et al.}
\newblock A \textit{Drosophila} computational brain model reveals sensorimotor
  processing.
\newblock \emph{Nature}, 634:210--219, 2024.

\bibitem{siegelmann1995}
H.~T. Siegelmann y E.~D. Sontag.
\newblock On the computational power of neural nets.
\newblock \emph{Journal of Computer and System Sciences}, 50(1):132--150, 1995.

\end{thebibliography}

%% ============================================================================
\appendix

\section{Programas de ejemplo}
\label{app:programas}

\begin{lstlisting}[style=dros,caption={\texttt{efe.dros}: emisión del carácter F.},label={lst:efe}]
;; efe.dros -- la primera palabra de la mosca
%cepa Canton-S

~~~> (DA1)                 ;; soplo de feromona sobre el glomérulo DA1
(DA1) !'F' x_x             ;; emitir 0x46 y morir
\end{lstlisting}

\begin{lstlisting}[style=dros,caption={\texttt{hola.dros}: saludo en UTF-8 mediante azúcar de cadena.},label={lst:hola}]
;; hola.dros -- saludo completo
%cepa w1118
~~~> (DM1)
(DM1) !"¡Hola, Mundo!\n" ~> (DA1)
(DA1) x_x
\end{lstlisting}

\begin{lstlisting}[style=dros,caption={\texttt{suma.dros}: suma por transferencia sináptica.},label={lst:suma}]
;; suma.dros -- la mosca que aprende a contar
%cepa Oregon-R

$ alfa = 0                         ;; sinapsis KC -> MBON, virgen
$ beta = 0                         ;; sinapsis KC -> MBON, virgen

~~~> (DM1)
(DM1) <alfa <beta ~> (VA2)         ;; percibir dos olores, grabar dos engramas
(VA2) ?beta ~> (VA3) | (DL5)       ;; ¿queda recuerdo en beta?
(VA3) +alfa ~> (VA2)               ;; pasar una unidad de beta a alfa
(DL5) !"suma: " ?alfa ~> (DA2) | (VM5d)   ;; primera lectura: si no hay nada, decir 0
(DA1) ?alfa ~> (DA2) | (VM7d)      ;; vaciar alfa por la probóscide
(DA2) !* ~> (DA1)                  ;; una extensión de probóscide por unidad
(VM5d) !'0' ~> (VM7d)              ;; suma nula: un 0 explícito
(VM7d) !'\n' x_x
\end{lstlisting}

\begin{lstlisting}[style=dros,caption={\texttt{longitud.dros}: longitud en bytes de un texto olido (\texttt{--olor-texto}).},label={lst:longitud}]
;; longitud.dros -- la mosca cuenta las letras de un texto olido
;; Cada byte llega como una concentración; el fin del texto huele a 0.
%cepa Canton-S

$ n = 0                            ;; engrama contador de bytes

~~~> (DM1)
(DM1) <c ~> (VA2)                  ;; oler el siguiente byte
(VA2) ?c ~> (VA3) | (DL5)          ;; olor nulo: se acabó el texto
(VA3) ?c ~> (VA3) | (VA4)          ;; olvidar el resto de ese olor
(VA4) +n ~> (DM1)                  ;; recordar un byte más
(DL5) !"bytes: " ?n ~> (DA2) | (VM5d)
(DA1) ?n ~> (DA2) | (VM7d)         ;; vaciar el contador por la probóscide
(DA2) !* ~> (DA1)
(VM5d) !'0' ~> (VM7d)
(VM7d) !'\n' x_x
\end{lstlisting}

\section{Trazas de ejecución}
\label{app:traza}

El \cref{lst:traza-efe} reproduce la traza completa de \texttt{efe.dros} y el
\cref{lst:traza-suma}, un extracto de \texttt{suma.dros} con entrada $(2,1)$.
Cada línea indica el paso, las neuronas reclutadas activas (papel y, entre
ángulos, tipo celular real), los pesos de los engramas y la actividad natural
por clase. Los símbolos no ASCII de la salida original se han transcrito.

\begin{lstlisting}[caption={Traza completa de \texttt{efe.dros}.},label={lst:traza-efe}]
t=0 | I(0): soplo sobre (DA1) = DA1_lPN . ALPN . right . id 720575940604407468
    | cerebro FlyWire v783: 138639 neuronas, 2700513 conexiones | implante: 13 neuronas
t=1 | DA1.0[!'F']<DA1_lPN>                                  | cerebro: 0
t=2 | bit1<CB0845> . bit2<CB0918> . bit6<CB0783> . estrobo<DNa02> . DA1.1[x_x]<KCg-m>
    | cerebro: 64 (Kenyon_Cell 48, central 14, ALPN 2)
t=3 | muerte<DNp01>                                         | cerebro: 12 (Kenyon_Cell 6, central 4, DAN 2)
+ FlyWire v783 . t = 3 . disparo de la Fibra Gigante . 76 neuronas despertaron
\end{lstlisting}

\begin{lstlisting}[caption={Extracto de la traza de \texttt{suma.dros} con entrada $(2,1)$.},label={lst:traza-suma}]
t=1  | DM1.0[<alfa]<DM1_lPN>               | alfa=0 beta=0 | cerebro: 0
t=2  | -                                    | alfa=0 beta=0 | cerebro: 462 (Kenyon_Cell 346, ...)
t=3  | ORN.alfa<ORN_DM2>                   | alfa=0 beta=0 | cerebro: 101
t=4  | PAM.alfa<PAM08> . ORN.alfa<ORN_DM2> | alfa=0 beta=0 | cerebro: 147
t=5  | PAM.alfa<PAM08>                     | alfa=1 beta=0 | cerebro: 78
t=6  | DM1.0.retorno<lLN2T_b>              | alfa=2 beta=0 | cerebro: 23
t=11 | DM1.1.retorno<CB1686>               | alfa=2 beta=1 | cerebro: 235
t=12 | VA2.0[?beta]<VA2_adPN>              | alfa=2 beta=1 | cerebro: 175
t=13 | KC.beta<KCg-s2> . VA2.retardo1      | alfa=2 beta=1 | cerebro: 500
t=14 | MBON.beta<MBON05> . VA2.retardo2    | alfa=2 beta=1 | cerebro: 151
t=15 | VA2.si->(VA3)<LHAV2m1>              | alfa=2 beta=1 | cerebro: 215
t=16 | PPL1.beta<PPL101> . VA3.0[+alfa]    | alfa=2 beta=1 | cerebro: 236
t=17 | PAM.alfa<PAM08> . VA2.0[?beta]      | alfa=2 beta=0 | cerebro: 300
t=18 | KC.beta<KCg-s2> . VA2.retardo1      | alfa=3 beta=0 | cerebro: 486
t=19 | VA2.retardo2<CB1058>                | alfa=3 beta=0 | cerebro: 203
t=20 | VA2.no->(DL5)<CB2211>               | alfa=3 beta=0 | cerebro: 231
...
t=47 | bit1 . bit3 . estrobo<DNa02> . VM7d.1[x_x]<KCab>   | alfa=0 beta=0 | cerebro: 265
t=48 | muerte<DNp01>                       | alfa=0 beta=0 | cerebro: 139
\end{lstlisting}

\end{document}
